\section{Estabilidad, regularización y programación}

\subsection{Entropy y regularización de la diversidad}

\begin{equation*}
\mathcal{L}_\pi = -\mathbb{E}_{s \sim \mathcal{D}, a \sim \pi} \left[\hat{A}(s,a)\log \pi(a|s)\right] + \beta \cdot \mathcal{H}(\pi(\cdot|s))
\end{equation*}

El término de entropy alienta a que la política conserve cierta aleatoriedad durante la etapa de actualización, con lo cual se evita que converja rápidamente hacia acciones subóptimas. Esto es especialmente importante en recompensas dispersas o en espacios de acción multimodales.

\begin{importantbox}{La magnitud práctica de entropy regularization}
En PPO / SAC es práctica común fijar el entropy bonus $\beta$ entre 0.01 y 0.1 de la loss de la política, es decir, conservar entre 1\% y 10\% de la energía de exploration. Si es demasiado pequeño se cae en un óptimo local; si es demasiado grande, se frena la convergencia de la política.
\end{importantbox}

\subsection{Trust region y clipped objective}

El clipped objective de PPO:
$${L}^{CLIP} = \mathbb{E} \left[ \min\left( r_t(\theta)\hat{A}_t, \; \text{clip}(r_t(\theta), 1-\epsilon, 1+\epsilon)\hat{A}_t \right) \right]$$
donde $r_t(\theta) = \frac{\pi_\theta(a_t|s_t)}{\pi_{\theta_{\text{old}}}(a_t|s_t)}$.

Tratar la restricción de KL o el mecanismo de clip como un trust region evita que la política dé pasos demasiado grandes. El docente enfatiza: en PPO hay que monitorear $KL(\pi_{\theta_{\text{old}}}||\pi_\theta)$ y, en cuanto se supere el umbral, regresar al checkpoint anterior.

\begin{knowledgebox}{Trust region contra clipping}
TRPO calcula con precisión la restricción sobre el gradiente, lo que exige mucho cómputo; PPO, en cambio, aproxima la restricción con clip, lo cual es sencillo de calcular y fácil de paralelizar, y es hoy el algoritmo dominante de elección.
\end{knowledgebox}

\subsection{Programación y target network}

El Target Network del critic (como los pesos lentos de SAC) suaviza el valor objetivo mediante Polyak averaging. Por ejemplo:
$$\theta_{\text{target}} \leftarrow \tau \theta + (1 - \tau) \theta_{\text{target}}$$

Un valor de $\tau$ más pequeño hace que el target network se actualice más lentamente, lo que mejora la estabilidad, pero también puede volverlo lento para reaccionar en entornos no estacionarios.

\begin{warningbox}{Si el target network se actualiza demasiado lento, no puede seguir el ritmo de la política}
Si $\tau$ se fija demasiado pequeño, el critic seguirá persiguiendo un target ya «desactualizado», lo cual hace que la estimación que guía al Actor sea demasiado optimista y que la adaptación a muestras nuevas sea lenta. Normalmente se fija $\tau$ entre 0.005 y 0.05.
\end{warningbox}

\subsection{Resumen de la sección}

Las técnicas de estabilidad (entropy, clip objective, target network) son decisivas para que Actor-Critic pueda converger con rapidez. Monitorear la distribución de KL y de advantage permite detectar a tiempo actualizaciones que se salen de rango.

